% LaTeX documentation converted from README.md
\documentclass[11pt]{article}
\usepackage[a4paper,margin=1in]{geometry}
\usepackage[hidelinks]{hyperref}
\usepackage{longtable}
\usepackage{enumitem}
\usepackage{fancyvrb}
\usepackage{titling}

\title{Bhadli Kavya -- Project Documentation}
\author{Bhadli Kavya Team}
\date{\today}

\begin{document}
\maketitle
\thispagestyle{empty}

\section*{Overview}
\begin{itemize}[leftmargin=*]
	\item \textbf{Name:} Bhadli Kavya -- Poetry \& Calendar Application
	\item \textbf{Stack:} React Native (Expo Web) frontend, FastAPI backend
	\item \textbf{Database:} Aiven Cloud MySQL (recommended) or local MySQL
	\item \textbf{Delivery:} Containerized with Docker and Docker Compose
\end{itemize}

\section{Quick Start for Developers}
\subsection{Prerequisites}
Install the following before attempting to run the project:
\begin{itemize}
	\item Docker -- see: \url{https://docs.docker.com/get-docker/}
	\item Docker Compose -- see: \url{https://docs.docker.com/compose/install/}
	\item Git
\end{itemize}

\section{Setup Instructions}
This section describes the steps to obtain and configure the project for development.

\subsection{1. Clone the Repository}
Open a terminal and run:
\begin{Verbatim}[fontsize=\small]
git clone https://github.com/erandesamadhan2003/bhadli-kavya.git
cd bhadli-kavya
\end{Verbatim}

\subsection{2. Configure Environment Variables}
Create a working \texttt{.env} from the template:
\begin{Verbatim}[fontsize=\small]
cp .env.example .env
\end{Verbatim}

Edit the new \texttt{.env} file (examples below) and fill in credentials and host information:
\begin{Verbatim}[fontsize=\small]
# Example variables to update:
DB_HOST=your-database-host
DB_PORT=your-database-port
DB_USER=your-database-user
DB_PASSWORD=your-database-password
DB_NAME=your-database-name
\end{Verbatim}

Important notes:
\begin{itemize}
	\item Use your own Aiven MySQL credentials or a local MySQL instance as described in the Database Configuration section.
	\item Do not commit the \texttt{.env} file to version control.
\end{itemize}

\subsection{3. Start the Application}
Use the provided script to build and start all services:
\begin{Verbatim}[fontsize=\small]
./start-docker.sh
\end{Verbatim}

This script performs the following actions:
\begin{itemize}
	\item Builds Docker images for frontend and backend
	\item Starts all services defined in \texttt{docker-compose.yml}
	\item Prints access URLs when startup completes
\end{itemize}

\section{Access the Application}
After containers start, use a browser to open the following URLs (default configuration):
\begin{longtable}{p{0.25\textwidth} p{0.55\textwidth}}
\textbf{Service} & \textbf{URL and Description} \\
\hline
Frontend & http://localhost:8081 -- React/Expo Web application \\
Backend API & http://localhost:8000 -- FastAPI REST API server \\
API Docs & http://localhost:8000/docs -- Swagger / interactive API documentation \\
\end{longtable}

\section{Development Commands}
Common docker-compose and maintenance commands used during development:
\subsection*{View Logs}
\begin{Verbatim}[fontsize=\small]
# All services
docker-compose logs -f

# Backend only
docker-compose logs -f backend

# Frontend only
docker-compose logs -f frontend
\end{Verbatim}

\subsection*{Stop / Restart / Status}
\begin{Verbatim}[fontsize=\small]
docker-compose stop
docker-compose restart
docker-compose ps
\end{Verbatim}

\section{Cleanup (When Done Working)}
To remove project containers and images and free disk space run:
\begin{Verbatim}[fontsize=\small]
./cleanup-docker.sh
\end{Verbatim}

What this script does:
\begin{itemize}
	\item Removes all project containers
	\item Removes all project images created for the project
	\item Frees up disk space
	\item Preserves your code files (no source deletion)
\end{itemize}

\section{Complete Developer Workflow}
Recommended daily workflow:
\begin{enumerate}
	\item Start services: \texttt{./start-docker.sh}
	\item Develop: frontend in \texttt{./frontend/}, backend in \texttt{./backend/}
	\item Monitor logs as needed: \texttt{docker-compose logs -f}
	\item When finished: \texttt{./cleanup-docker.sh}
	\item Next session: \texttt{./start-docker.sh}
\end{enumerate}

\section{Project Structure}
The repository layout (top-level) is:
\begin{Verbatim}[fontsize=\small]
bhadli-kavya/
├── backend/              # FastAPI backend
│   ├── Dockerfile
│   ├── main.py
│   ├── database.py
│   ├── models.py
│   ├── routers/
│   └── requirements.txt
├── frontend/             # React/Expo frontend
│   ├── Dockerfile
│   ├── App.jsx
│   ├── package.json
│   ├── components/
│   ├── screens/
│   └── services/
├── docker-compose.yml    # Service orchestration
├── start-docker.sh       # Start script
├── cleanup-docker.sh     # Cleanup script
├── .env                  # Your environment variables (DO NOT COMMIT)
└── .env.example          # Environment template
\end{Verbatim}

\section{Database Configuration}
Two supported options are provided below.

\subsection{Option 1: Aiven Cloud MySQL (Recommended)}
Steps:
\begin{enumerate}
	\item Create a MySQL database on \url{https://aiven.io} or your cloud provider
	\item Obtain connection credentials from the Aiven console
	\item Update the project's \texttt{.env} file with the Aiven values
\end{enumerate}

Example Aiven configuration:
\begin{Verbatim}[fontsize=\small]
DB_HOST=your-aiven-host.aivencloud.com
DB_PORT=21638
DB_USER=avnadmin
DB_PASSWORD=your-aiven-password
DB_NAME=defaultdb
DB_SSL_MODE=REQUIRED
\end{Verbatim}

\subsection{Option 2: Local MySQL Database}
To use a local database instead of a managed Aiven instance:
\begin{enumerate}
	\item Uncomment the \texttt{database} service in \texttt{docker-compose.yml}
	\item Update the \texttt{.env} file with the local database credentials
	\item Example variables:
\end{enumerate}

\begin{Verbatim}[fontsize=\small]
DB_HOST=database
DB_PORT=3306
DB_USER=your_username
DB_PASSWORD=your_password
DB_NAME=bhadli_kavya
DB_SSL_MODE=DISABLED
\end{Verbatim}

\section{Troubleshooting}
Common problems and solutions.

\subsection{Port Already in Use}
If ports 8000 or 8081 conflict with other services:
\begin{enumerate}
	\item Edit the \texttt{.env} file to change ports:
\begin{Verbatim}[fontsize=\small]
API_PORT=9000
FRONTEND_PORT=9001
\end{Verbatim}
	\item Restart with \texttt{./start-docker.sh}
\end{enumerate}

\subsection{Database Connection Error}
Verify your \texttt{.env} credentials and then inspect backend logs:
\begin{Verbatim}[fontsize=\small]
cat .env
docker-compose logs backend
\end{Verbatim}

\subsection{Container Won't Start}
Rebuild and start fresh:
\begin{Verbatim}[fontsize=\small]
./cleanup-docker.sh
./start-docker.sh
\end{Verbatim}

\section{Contributing}
To contribute to the project:
\begin{enumerate}
	\item Fork the repository
	\item Create a feature branch: \texttt{git checkout -b feature/amazing-feature}
	\item Commit: \texttt{git commit -m 'Add amazing feature'}
	\item Push: \texttt{git push origin feature/amazing-feature}
	\item Open a pull request for review
\end{enumerate}

\section{Important Notes}
\begin{itemize}
	\item \textbf{Security:} Never commit the \texttt{.env} file. It contains sensitive credentials.
	\item \textbf{Template:} Use \texttt{.env.example} as the template for environment variables.
	\item \textbf{Hot reload:} Containers are configured to hot-reload source changes where supported; code changes should reflect immediately.
	\item \textbf{Persistence:} Database data persists between restarts when using the local database configuration.
\end{itemize}

\section{Getting Help}
If you run into issues:
\begin{itemize}
	\item Check logs: \texttt{docker-compose logs -f}
	\item Check status: \texttt{docker-compose ps}
	\item Rebuild everything if necessary: \texttt{./cleanup-docker.sh \&\& ./start-docker.sh}
\end{itemize}

\section*{Final Remarks}
To quickly get started run:
\begin{Verbatim}[fontsize=\small]
./start-docker.sh
# then open: http://localhost:8081
\end{Verbatim}

\vfill
\noindent{\small Made with love by the Bhadli Kavya Team}

\end{document}

